\subsection{Theorem A: proof and verification method}\label{sec:methodA}
Theorem~\ref{thm:Ainterval} is proved directly from the game rules. A move $r$ from $(n,s)$ reaches a losing state exactly when $r\in A_n$ and $|r-s|>m$, so $(n,s)$ is losing precisely when every $r\in A_n$ satisfies $|r-s|\le m$. The admissible $s$ are therefore the intersection of the intervals $[r-m,r+m]$ over $r\in A_n$, clipped to $[1,k]$: a single interval of at most $2m+1$ entries, or the full row when $A_n=\varnothing$. Finite memory holds because $A_n$ uses only $Q_{n-1},\ldots,Q_{n-k}$ for $n\ge k$; eventual periodicity follows by pigeonhole on the finitely many $k$-row windows, and the starting-position predicate inherits it.
The computation is verified, not trusted: the interval recurrence is checked against direct legal-move quantification and an independent Sprague--Grundy solver for $k=1,\ldots,12$ and $m=0,\ldots,4$ (the \texttt{small\_checks} routine in \texttt{methodology/theorem\_B/verify\_theorems\_ab.py}), alongside the solver self-tests in \texttt{src/}.
